\documentclass{filkomskripsi}
% Ganti nilai contoh ini dengan identitas dan judul yang telah disetujui.
% Kosongkan font-dir bila menggunakan font yang terpasang pada sistem.
% Berkas font berlisensi dapat disiapkan sendiri mengikuti petunjuk proyek.
\filkomsetup{
  judul={Contoh Judul Skripsi yang Menjelaskan Fokus Penelitian},
  judul-inggris={A Sample Thesis Title that Describes the Research Focus},
  nama={Nama Mahasiswa},
  nim={000000000000000},
  program-studi={TEKNIK INFORMATIKA},
  departemen={DEPARTEMEN TEKNIK INFORMATIKA},
  fakultas={FAKULTAS ILMU KOMPUTER},
  universitas={UNIVERSITAS BRAWIJAYA},
  kota={MALANG},
  tahun={2026},
  pembimbing-1={Nama Pembimbing 1},
  pembimbing-2={Nama Pembimbing 2},
  font-dir={}
}

\makeatletter
\def\input@path{{bibliography/}}
\makeatother
\usepackage[backend=biber,style=filkom-authoryear,language=english,autolang=none]{biblatex}
\DeclareLanguageMapping{indonesian}{english}
\ExecuteBibliographyOptions{uniquename=false,uniquelist=false,urldate=iso}
% Tanggal akses bukan tahun penerbitan; tanpa date/year tampil n.d.
\DeclareLabeldate{\field{date}\field{year}\literal{nodate}}
\setlength{\bibhang}{0.9cm}
\addbibresource{bibliography/references.bib}
\usepackage{tikz}
\usetikzlibrary{arrows.meta,positioning,shapes.geometric}
\emergencystretch=3em
\begin{document}
\makeskripsicover
\skripsifrontmatter
\skripsifrontchapter{PERSETUJUAN}

\begin{center}
  \textbf{\judulskripsi}\par
  \vspace{0.7cm}
  \namamahasiswa\par
  NIM: \nimmahasiswa\par
\end{center}

\vspace{0.6cm}
Halaman ini merupakan contoh formulir kosong. Redaksi persetujuan, keperluan pengajuan, dan pihak yang mengesahkan disesuaikan dengan formulir fakultas yang berlaku. Dokumen contoh belum menyatakan adanya persetujuan, kelulusan, atau pengesahan akademik.

\vspace{0.5cm}
Keperluan persetujuan: \rule{7.0cm}{0.4pt}\par
Tanggal persetujuan: \rule{7.1cm}{0.4pt}\par

\vspace{1.0cm}
\begin{center}
  \begin{tabular}{p{5.8cm}p{5.8cm}}
    Dosen Pembimbing 1 & Dosen Pembimbing 2 \\
    \rule{0pt}{2.5cm} & \\
    \pembimbingsatu & \pembimbingdua \\
    NIP/NIK: \rule{3.4cm}{0.4pt} & NIP/NIK: \rule{3.4cm}{0.4pt} \\
  \end{tabular}
\end{center}

% Isi tanggal dan tanda tangan hanya berdasarkan persetujuan yang sebenarnya.

% Setelah ujian dan revisi disahkan, ganti input di atas dengan:
% % Gunakan setelah ujian dan revisi benar-benar disahkan.
% Lengkapi nama penguji, identitas pejabat, tanggal, serta tanda tangan aktual.
\skripsifrontchapter{PENGESAHAN}
\begin{center}
\textbf{\judulskripsi}\par
\namamahasiswa\par NIM: \nimmahasiswa
\end{center}
Naskah ini dipertahankan di hadapan penguji pada tanggal \rule{3cm}{0.4pt}. Hasil ujian dan penyelesaian revisi diisi sesuai berita acara yang telah disahkan.

\begin{center}
\begin{tabular}{p{5.8cm}p{5.8cm}}
Pembimbing 1 & Pembimbing 2 \\
\rule{0pt}{2cm} & \\
\pembimbingsatu & \pembimbingdua \\
\end{tabular}
\end{center}
\vspace{0.8cm}
\noindent Penguji: \rule{7cm}{0.4pt}\par
\noindent Ketua Departemen: \rule{7cm}{0.4pt}\par
\noindent Tanggal pengesahan: \rule{5cm}{0.4pt}

\skripsifrontchapter{PERNYATAAN ORISINALITAS}

\begin{tabular}{lp{9.5cm}}
  Nama & \namamahasiswa \\
  NIM & \nimmahasiswa \\
  Judul & \judulskripsi \\
\end{tabular}

\vspace{0.7cm}
Halaman ini menyediakan ruang untuk pernyataan orisinalitas sesuai formulir yang berlaku. Masukkan redaksi resmi yang dapat dipertanggungjawabkan setelah naskah disusun dan diperiksa. Materi pada dokumen ini merupakan contoh format; halaman ini belum menjadi pernyataan yang ditandatangani.

\vspace{0.5cm}
\noindent Redaksi pernyataan:\par
\noindent\rule{\textwidth}{0.4pt}\par
\noindent\rule{\textwidth}{0.4pt}\par
\noindent\rule{\textwidth}{0.4pt}\par
\noindent\rule{\textwidth}{0.4pt}\par

\vspace{1.2cm}
\begin{flushright}
  Tempat, tanggal: \rule{4.0cm}{0.4pt}\par
  \vspace{2.3cm}
  \namamahasiswa\par
  NIM: \nimmahasiswa
\end{flushright}

\skripsifrontchapter{PRAKATA}

Prakata pada dokumen ini merupakan contoh petunjuk penulisan. Pengguna dapat menjelaskan konteks penyusunan skripsi, menyampaikan terima kasih atas bantuan yang benar-benar diterima, serta menuliskan harapan yang sesuai dengan isi karya. Uraian tidak perlu mengulang abstrak, latar belakang, atau simpulan penelitian.

Pihak yang disebut dalam ucapan terima kasih dipilih berdasarkan kontribusinya. Nama, gelar, jabatan, dan bentuk bantuan diperiksa agar penulis tidak memberikan keterangan yang keliru. Bantuan akademik, dukungan administratif, serta dukungan pribadi dapat dijelaskan secara singkat tanpa mengungkap informasi yang tidak diperlukan dalam naskah publik.

Sebelum naskah digunakan, ganti seluruh paragraf petunjuk ini dengan prakata penulis. Tempat dan tanggal pada bagian akhir diisi sesuai penyelesaian atau ketentuan penyerahan yang berlaku. Ruang di bawah sengaja tidak memuat tanggal ataupun identitas pribadi nyata.

\vspace{1.5cm}
\begin{flushright}
  Tempat, tanggal: \rule{4.0cm}{0.4pt}\par
  \vspace{1.6cm}
  \namamahasiswa
\end{flushright}

\skripsifrontchapter{ABSTRAK}

\noindent\textbf{\namamahasiswa, \judulskripsi}\par
\noindent\textbf{Pembimbing: \pembimbingsatu{} dan \pembimbingdua}\par
\vspace{\baselineskip}

Naskah contoh ini memperlihatkan penggunaan template skripsi berbasis LaTeX untuk menata bagian awal, enam bab utama, daftar referensi, dan lampiran. Tujuannya adalah menyediakan titik awal yang dapat disunting serta membantu pengguna mengenali hubungan antara struktur sumber dan tampilan dokumen. Materi yang disertakan merupakan penjelasan format, bukan laporan penelitian empiris. Identitas mahasiswa dan pembimbing menggunakan penanda generik, sedangkan ruang persetujuan, pernyataan orisinalitas, tanggal, dan tanda tangan tetap memerlukan pengisian berdasarkan keadaan sebenarnya. Bab pendahuluan menunjukkan cara menghubungkan konteks, pertanyaan, tujuan, manfaat, dan batasan. Bab landasan memberikan contoh sitasi dokumentasi resmi serta definisi simbol matematika. Bab metodologi menampilkan hubungan tahap dan bahan pemeriksaan melalui tabel dan diagram vektor yang dibuat langsung di dalam naskah. Bab hasil memperlihatkan tabel ringkas, tabel yang melintasi halaman dengan judul kolom berulang, serta rujukan silang ke gambar dan persamaan. Isi tabel menggunakan nama komponen dan deskripsi keluaran tanpa skor, jumlah sampel, atau keberhasilan pengujian yang direka. Bab pembahasan menerangkan hubungan bukti, penafsiran, dan keterbatasan, sementara penutup menunjukkan bentuk simpulan serta saran yang perlu disesuaikan. Daftar referensi berisi sumber daring resmi untuk memperlihatkan pemanggilan bibliografi, dan lampiran menata catatan pendukung. Pengguna perlu mengganti seluruh materi contoh dengan isi penelitian yang dapat dipertanggungjawabkan, memeriksa aturan penulisan yang berlaku, serta membaca kembali keluaran setelah perubahan terakhir. Dokumen yang berhasil dikompilasi tetap memerlukan pemeriksaan isi, format, rujukan, dan kelengkapan administratif sebelum digunakan sebagai naskah akademik.

\noindent\textbf{Kata kunci:} LaTeX, template skripsi, struktur naskah, sitasi, rujukan silang, pemeriksaan format.

\skripsifrontchapter{ABSTRACT}

\begin{otherlanguage}{english}
\noindent\textbf{\namamahasiswa, \judulskripsiinggris}\par
\noindent\textbf{Supervisors: \pembimbingsatu{} and \pembimbingdua}\par
\vspace{\baselineskip}

This sample document demonstrates a LaTeX thesis template for organizing preliminary pages, six main chapters, a bibliography, and an appendix. Its purpose is to provide an editable starting point and help users understand the relationship between source structure and document appearance. The included material explains formatting and does not report an empirical study. Student and supervisor identities use generic values, while approval statements, declarations of originality, dates, and signatures require completion according to actual circumstances. The introduction illustrates connections among context, questions, objectives, benefits, and scope. The background chapter provides examples of citations to official documentation and definitions of mathematical symbols. The methodology chapter presents relationships between stages and materials through a table and a vector diagram created within the document. The results chapter demonstrates a compact table, a multipage table with repeated column headings, and cross references to a figure and equations. Its tables contain component names and output descriptions without invented scores, sample sizes, or successful experiments. The discussion chapter explains relationships among evidence, interpretation, and limitations, while the closing chapter illustrates conclusions and recommendations that users must adapt. The bibliography contains genuine online documentation to demonstrate citation management, and the appendix organizes supporting notes. Users must replace every sample passage with accountable research content, check the applicable writing requirements, and review the generated document after the final changes. Successful compilation does not replace examination of content, formatting, references, or administrative completeness before the document is used as an academic submission.

\noindent\textbf{Keywords:} LaTeX, thesis template, document structure, citation, cross reference, format review.
\end{otherlanguage}

\makeskripsilists
\skripsimainmatter
\chapter{PENDAHULUAN}
\label{bab:pendahuluan}

Bab ini memperlihatkan susunan pendahuluan pada contoh naskah skripsi. Seluruh uraian merupakan petunjuk penulisan dan contoh format yang perlu diganti dengan materi penelitian pengguna. Naskah ini tidak mewakili proposal yang telah disetujui, penelitian yang telah dilaksanakan, ataupun hasil pengujian suatu metode. Tujuan contoh adalah memperlihatkan hubungan antara bagian naskah, penomoran, dan penyajian unsur akademik dalam satu dokumen yang dapat dikompilasi.

\section{Latar Belakang}

Pendahuluan membantu pembaca mengenali masalah, konteks, serta alasan suatu penelitian perlu dilakukan. Paragraf awal dapat memperkenalkan keadaan umum yang relevan dengan objek penelitian. Paragraf berikutnya mempersempit konteks hingga persoalan yang dapat diamati dan diteliti. Alur tersebut memungkinkan pembaca memahami hubungan antara kebutuhan penelitian dan pendekatan yang dipilih, tanpa harus menebak latar pemilihan topik.

Dalam naskah penelitian sesungguhnya, uraian latar belakang perlu ditopang sumber yang telah dibaca dan diperiksa. Pernyataan tentang keterbatasan pendekatan terdahulu harus sesuai dengan isi sumber yang dirujuk. Jika keterbatasan hanya berlaku pada objek atau kondisi tertentu, batas tersebut perlu dipertahankan dalam kalimat. Sitasi tidak menggantikan penjelasan mengenai bagaimana sumber mendukung argumentasi penulis.

Contoh naskah ini menggunakan penyusunan dokumen sebagai konteks pembelajaran format. Penulis memisahkan isi naskah dari pengaturan tampilannya, sehingga perubahan judul bagian atau penambahan tabel dapat diikuti oleh pembaruan nomor dan rujukan. Pada penelitian sesungguhnya, konteks ini diganti dengan persoalan substantif yang menjadi fokus skripsi. Istilah teknis yang belum dikenali pembaca perlu dijelaskan sebelum digunakan untuk membangun argumen.

Penutup latar belakang sebaiknya menghubungkan masalah dengan rencana jawaban penelitian. Hubungan itu diperjelas melalui pertanyaan, tujuan, serta batasan yang saling sesuai. Bab ini menunjukkan bentuk hubungan tersebut; Bab \ref{bab:metodologi} memperlihatkan cara menata tahapan, sedangkan Bab \ref{bab:hasil} menunjukkan tempat penyajian keluaran. Rujukan silang tersebut akan mengikuti perubahan nomor bab secara otomatis.

\section{Rumusan Masalah}

Rumusan masalah ditulis sebagai pertanyaan yang dapat dijawab oleh prosedur dan bukti penelitian. Pada contoh penulisan ini, pertanyaannya digunakan untuk mengarahkan penggunaan template:

\begin{enumerate}
  \item Bagaimana menyusun naskah dengan hubungan yang jelas antara pendahuluan, landasan, metodologi, hasil, pembahasan, dan penutup?
  \item Bagaimana menyajikan tabel, gambar, persamaan, sitasi, dan lampiran yang mempunyai rujukan silang konsisten?
  \item Bagaimana memeriksa kesesuaian keluaran dokumen dengan sumber yang dapat disunting kembali?
\end{enumerate}

Pengguna perlu mengganti ketiga pertanyaan tersebut dengan rumusan masalah penelitiannya. Jumlah pertanyaan tidak perlu mengikuti jumlah pada contoh. Setiap pertanyaan harus memiliki kebutuhan bukti yang jelas; pertanyaan yang hanya menjelaskan kegiatan pengerjaan dapat dipertimbangkan kembali apabila tidak menghasilkan jawaban ilmiah.

\section{Tujuan}

Tujuan penulisan contoh ini adalah menyediakan susunan naskah yang dapat dipelajari dan disesuaikan. Keterkaitannya dengan rumusan masalah diperlihatkan pada Tabel \ref{tab:contoh-keterkaitan}. Tabel ini tidak menunjukkan capaian penelitian. Pengguna dapat memakai pola yang sama untuk memastikan setiap tujuan memiliki cara pemenuhan dan tempat pelaporan yang sesuai.

\begin{table}[htbp]
  \centering
  \caption{Contoh keterkaitan pertanyaan, tujuan, dan bagian pelaporan}
  \label{tab:contoh-keterkaitan}
  \begin{tabular}{p{3.6cm}p{4.2cm}p{4.4cm}}
    \toprule
    Fokus pertanyaan & Tujuan penulisan contoh & Bagian pelaporan \\
    \midrule
    Susunan naskah & Menunjukkan hubungan antarbagian & Pendahuluan sampai penutup \\
    Unsur dokumen & Menunjukkan rujukan silang & Tabel, gambar, persamaan, dan lampiran \\
    Pemeriksaan keluaran & Menjelaskan cara membaca keluaran & Pembahasan dan lampiran \\
    \bottomrule
  \end{tabular}
\end{table}

\section{Manfaat}

Manfaat contoh naskah bagi penulis adalah tersedianya titik awal yang mempunyai struktur dan unsur dokumen yang dapat diamati. Penulis dapat mengganti materi sambil mempertahankan pola penomoran serta rujukan yang diperlukan. Bagi pembaca, contoh ini memperlihatkan hubungan antara suatu unsur dokumen dan penjelasan yang menyertainya. Manfaat tersebut bersifat pembelajaran tata tulis, bukan klaim manfaat dari suatu metode penelitian.

Pada skripsi sesungguhnya, manfaat perlu dijelaskan menurut pihak yang menerima kegunaan hasil. Manfaat teoretis dapat dikaitkan dengan pengetahuan yang diperoleh, sedangkan manfaat praktis dapat dikaitkan dengan keputusan atau kegiatan yang terbantu. Pernyataan manfaat tidak sebaiknya menjanjikan sesuatu di luar ruang lingkup dan bukti penelitian.

\section{Batasan}

Contoh ini memperlihatkan struktur enam bab untuk susunan nonimplementatif. Susunan tersebut dapat disesuaikan dengan jenis penelitian, pedoman yang berlaku, dan arahan pembimbing. Nama bagian serta redaksi administratif perlu diperiksa kembali sebelum naskah digunakan sebagai dokumen akademik. Persetujuan atas isi, format akhir, dan kelengkapan administratif tidak timbul hanya karena dokumen dapat dikompilasi.

Isi contoh dibatasi pada penataan naskah. Tidak ada dataset, partisipan, skor evaluasi, ataupun hasil eksperimen yang disertakan. Tabel hasil menggunakan nama komponen dan deskripsi keluaran, sehingga tidak menyajikan angka empiris yang dibuat untuk mengisi halaman. Contoh rumus merupakan penjelasan notasi matematika dan dapat diganti dengan rumus yang relevan terhadap penelitian pengguna.

\section{Sistematika Penulisan}

Bab \ref{bab:pendahuluan} memperkenalkan konteks, pertanyaan, tujuan, manfaat, serta batas contoh. Bab \ref{bab:landasan} menyajikan konsep penataan dokumen dan contoh sitasi sumber resmi. Bab \ref{bab:metodologi} menunjukkan susunan prosedur dan diagram alur. Bab \ref{bab:hasil} memperlihatkan format pelaporan tanpa mengklaim hasil empiris. Bab \ref{bab:pembahasan} menunjukkan hubungan bukti dan penafsiran, sedangkan Bab \ref{bab:penutup} menata simpulan serta saran. Materi pendukung ditempatkan pada Lampiran \ref{lamp:contoh}.

\chapter{LANDASAN KEPUSTAKAAN}
\label{bab:landasan}

Bab ini menunjukkan cara menata landasan, mendefinisikan istilah, serta menghubungkan uraian dengan sumber. Sumber yang digunakan merupakan dokumentasi resmi perangkat penyusunan naskah. Rujukan tersebut disertakan untuk memperlihatkan sitasi yang dapat ditelusuri, bukan sebagai pengganti tinjauan pustaka pada bidang penelitian pengguna.

\section{Landasan Penyusunan Dokumen}

\subsection{Struktur dan Makna Bagian}

Penataan dokumen dimulai dari hubungan antargagasan. Bab mengelompokkan pokok bahasan utama, bagian memperinci pokok tersebut, dan subbagian memisahkan uraian yang masih berada dalam satu pembahasan. Pembagian yang jelas membantu pembaca menemukan informasi tanpa menghilangkan alur argumentasi. Pemilihan judul sebaiknya didasarkan pada isi, bukan pada kebutuhan untuk mengisi setiap tingkat penomoran.

\textcite{latexproject-documentation} menyediakan dokumentasi inti \LaTeX{} bagi penulis, termasuk panduan tentang perintah dan penyusunan dokumen. Dalam contoh ini, perintah judul, rujukan, dan persamaan digunakan untuk menunjukkan pemisahan struktur dari tampilan. Satu perubahan pada pengaturan dokumen dapat diterapkan ke bagian terkait tanpa mengubah redaksi setiap paragraf secara manual.

\subsection{Paragraf dan Penjelasan Unsur}

Paragraf yang baik mengembangkan satu pokok gagasan. Kalimat pembuka dapat menyatakan pokok tersebut, kemudian kalimat berikutnya memberi penjelasan, bukti, atau hubungan dengan pembahasan. Paragraf penutup bagian dapat menghubungkan gagasan yang telah diterangkan dengan langkah berikutnya. Panjang paragraf disesuaikan dengan kebutuhan penjelasan, sehingga pembaca dapat mengikuti gagasan tanpa memecah atau menggabungkannya kembali.

Tabel dan gambar membantu menjelaskan informasi tertentu, tetapi tetap memerlukan pengantar serta penjelasan. Pengantar memberi alasan unsur tersebut ditampilkan. Penjelasan sesudah unsur dapat mengarahkan pembaca pada hubungan yang relevan dengan tujuan bagian. Menempatkan tabel tanpa merujuknya dalam narasi membuat pembaca harus menafsirkan kegunaannya sendiri.

\subsubsection{Contoh Tingkat Judul Ketiga}

Subsubbagian ini memperlihatkan tingkat judul ketiga. Gunakan tingkat ini apabila pembagian tambahan membantu pembaca mengenali uraian yang masih berada dalam satu subbagian. Apabila sebuah subbagian hanya memerlukan paragraf singkat, judul tambahan mungkin tidak diperlukan. Pengguna dapat menghapus bagian ini setelah memahami perilaku penomorannya.

\section{Font dan Mesin Kompilasi}

Paket \texttt{fontspec} menyediakan antarmuka pemilihan font untuk XeLaTeX dan LuaLaTeX, termasuk dukungan font OpenType \parencite{ctan-fontspec}. Informasi tersebut mendasari penggunaan mesin kompilasi yang sesuai dalam template. Ketersediaan font dan hak penggunaannya perlu diperhatikan ketika naskah dipindahkan antarperangkat atau dibagikan kepada pengguna lain.

Pilihan font pada keluaran perlu dibaca bersama ketentuan penulisan yang berlaku. Tampilan yang mirip belum tentu memenuhi ketentuan font tertentu. Ketika pengguna menyiapkan berkas font sendiri, pengaturan proyek sebaiknya dibuat jelas agar dokumen dapat diperiksa dengan lingkungan yang sama. Naskah contoh tidak menyertakan berkas font berlisensi milik pengguna.

\section{Sitasi dan Daftar Referensi}

\subsection{Sitasi dalam Narasi}

Sitasi dapat menjadi bagian dari narasi ketika penulis atau organisasi sumber perlu disebut secara langsung. Bentuk \textcite{ctan-biblatex} pada kalimat ini memperlihatkan penempatan tersebut. Bentuk lain meletakkan identitas sumber dalam tanda kurung pada akhir pernyataan yang didukung \parencite{ctan-biblatex}. Pengguna dapat membandingkan kedua bentuk pada sumber dan keluaran naskah.

Paket \texttt{biblatex} mengelola pemformatan sitasi dan daftar referensi, serta dapat memakai Biber untuk membaca data bibliografis \parencite{ctan-biblatex}. Dalam contoh ini, identitas sumber ditulis pada berkas bibliografi dan dipanggil melalui kunci sitasi. Jika nama sumber atau alamat berubah setelah pemeriksaan, pembaruan dilakukan pada entri sumber sehingga penulisan dalam naskah tetap konsisten.

\subsection{Kelengkapan Data Bibliografis}

Setiap entri harus sesuai dengan jenis sumber sebenarnya. Artikel jurnal, buku, dan halaman dokumentasi mempunyai unsur identitas yang berbeda. Penulis tidak perlu mengisi DOI, tahun, nomor volume, atau halaman dengan nilai rekaan hanya agar semua kolom tampak terisi. Apabila tahun penerbitan halaman daring tidak diketahui, data tersebut dapat tidak dicantumkan; tanggal akses dicatat secara terpisah.

Bibliografi contoh memuat tiga halaman dokumentasi resmi bertipe sumber daring. Semuanya dipanggil pada bab ini sehingga hubungan antara sitasi dan daftar referensi dapat diperiksa. Pada penelitian sesungguhnya, tambahkan hanya sumber yang digunakan dalam argumentasi. Daftar referensi tidak perlu diperluas dengan sumber yang belum dibaca untuk memperbesar jumlah entri.

\section{Contoh Notasi dan Persamaan}
\label{sec:contoh-notasi}

Persamaan \ref{eq:contoh-rerata} memperlihatkan penulisan rerata aritmetika untuk himpunan nilai $x_1,\ldots,x_n$. Persamaan tersebut hanya contoh format, tanpa nilai observasi yang dinyatakan telah diperoleh:
\begin{equation}
  \bar{x}=\frac{1}{n}\sum_{i=1}^{n}x_i,\qquad n>0.
  \label{eq:contoh-rerata}
\end{equation}
Indeks $i$ menunjuk anggota himpunan, $n$ menyatakan banyaknya anggota, dan $\bar{x}$ menyatakan rerata. Notasi didefinisikan setelah persamaan agar pembaca tidak perlu menebak arti simbol. Pada skripsi sesungguhnya, jelaskan asumsi, satuan, serta kondisi ketika suatu persamaan dapat digunakan.

Contoh lain adalah transformasi berdasarkan dua batas nilai yang berbeda. Persamaan \ref{eq:contoh-normalisasi} memperlihatkan pecahan, indeks, dan syarat penyebut:
\begin{align}
  z_i &= \frac{x_i-x_{\min}}{x_{\max}-x_{\min}},
  \label{eq:contoh-normalisasi}\\
  x_{\max} &> x_{\min}.
  \nonumber
\end{align}
Transformasi ini ditampilkan untuk mengilustrasikan lingkungan persamaan beberapa baris. Contoh tidak menyatakan bahwa transformasi tersebut selalu diperlukan atau tepat bagi penelitian tertentu. Definisi besaran dan batas penggunaannya harus disesuaikan dengan metode sebenarnya.

Tabel \ref{tab:contoh-notasi} menata simbol yang telah digunakan. Pada naskah dengan banyak simbol, tabel serupa dapat membantu menjaga kesamaan makna antarbab.

\begin{table}[htbp]
  \centering
  \caption{Contoh keterangan notasi matematika}
  \label{tab:contoh-notasi}
  \begin{tabular}{lp{9.8cm}}
    \toprule
    Simbol & Keterangan \\
    \midrule
    $x_i$ & Nilai anggota ke-$i$ dalam suatu himpunan \\
    $n$ & Banyaknya anggota, dengan syarat $n>0$ \\
    $\bar{x}$ & Rerata aritmetika sebagaimana Persamaan \ref{eq:contoh-rerata} \\
    $x_{\min},x_{\max}$ & Batas bawah dan batas atas yang berbeda \\
    $z_i$ & Nilai hasil transformasi pada Persamaan \ref{eq:contoh-normalisasi} \\
    \bottomrule
  \end{tabular}
\end{table}

\section{Hubungan Landasan dengan Metode}

Landasan kepustakaan sebaiknya berakhir dengan hubungan konsep dan keputusan penelitian. Konsep yang dijelaskan perlu mempunyai peran dalam pemilihan metode, penyusunan instrumen, pembentukan ukuran, atau penafsiran. Sebaliknya, konsep yang tidak digunakan dapat dipertimbangkan kembali agar pembaca tidak menerima uraian panjang yang terputus dari tujuan penelitian.

Pada contoh ini, struktur dokumen mendasari pembagian bab, rujukan silang menghubungkan unsur, dan data bibliografis mengikat pernyataan dengan sumbernya. Bab \ref{bab:metodologi} memakai hubungan tersebut untuk menata tahapan penyusunan contoh. Pengguna perlu mengganti hubungan ini dengan hubungan yang benar pada penelitian yang dikerjakan.

\chapter{METODOLOGI}
\label{bab:metodologi}

Bab ini menunjukkan bentuk uraian metodologi pada susunan enam bab. Prosedur contoh berfokus pada penyiapan dan pemeriksaan naskah. Tidak ada pengumpulan data atau eksperimen ilmiah yang dinyatakan telah dilaksanakan. Dalam skripsi pengguna, seluruh bagian berikut perlu diganti dengan prosedur yang memang digunakan untuk menjawab rumusan masalah.

\section{Rancangan}

Rancangan menerangkan hubungan antara tujuan, sumber informasi, prosedur, dan keluaran yang akan dianalisis. Penjelasan tersebut memungkinkan pembaca menilai apakah metode cukup untuk menghasilkan jawaban. Nama metode tidak cukup menjelaskan rancangan; pembaca juga memerlukan unit yang diamati, cara pemilihan unit, urutan perlakuan, serta aturan pemeriksaan.

Pada contoh penyusunan naskah, unit yang diperiksa adalah bagian dokumen, seperti judul, paragraf, tabel, gambar, persamaan, sitasi, dan lampiran. Setiap unit memiliki peran yang berbeda dalam penyampaian argumentasi. Prosedur diarahkan untuk memastikan peran tersebut terbaca dan hubungan antarbagiannya dapat ditelusuri melalui rujukan silang.

\section{Sumber dan Bahan}

Sumber pada naskah ini berupa dokumentasi resmi yang dirujuk dalam Bab \ref{bab:landasan}, sumber \LaTeX{} yang dapat disunting, serta keluaran dokumen yang perlu diperiksa. Ketiga kelompok tersebut mempunyai fungsi berbeda. Dokumentasi menerangkan kemampuan perangkat, sumber menyimpan isi dan penanda struktur, sedangkan keluaran memperlihatkan tampilan yang diterima pembaca.

Pada penelitian sesungguhnya, sumber dan bahan perlu dijelaskan menurut asal, cakupan, waktu, izin penggunaan, serta cara pemerolehan. Jika data diperoleh dari pihak lain, prosedur tidak boleh memberi kesan bahwa penulis mengumpulkannya secara langsung. Bila suatu bahan belum tersedia, kondisi itu perlu dilaporkan tanpa memperkirakan ukuran atau isinya.

\section{Tahapan}

Gambar \ref{fig:contoh-alur} menunjukkan urutan umum penyusunan contoh. Diagram dibuat dengan unsur vektor di dalam naskah sehingga tidak bergantung pada gambar dari berkas tambahan. Panah menjelaskan hubungan berurutan, bukan hubungan sebab-akibat ilmiah. Pada penelitian pengguna, nama tahap dan hubungan antarproses harus sesuai dengan metode yang benar-benar dikerjakan.

\begin{figure}[htbp]
  \centering
  \begin{tikzpicture}[
    every node/.style={font=\small,align=center},
    tahap/.style={draw,rounded corners=2pt,text width=9.8cm,
      minimum height=0.95cm,inner sep=5pt}
  ]
    \node[tahap] (arah) at (0,0) {Menentukan fokus, struktur, dan kebutuhan penjelasan};
    \node[tahap] (sumber) at (0,-1.5) {Menyiapkan sumber dan data bibliografis yang dapat diperiksa};
    \node[tahap] (isi) at (0,-3.0) {Menulis isi serta memberi penanda tabel, gambar, dan persamaan};
    \node[tahap] (keluaran) at (0,-4.5) {Menyusun keluaran dokumen dan memeriksa rujukan silang};
    \node[tahap] (perbaikan) at (0,-6.0) {Memperbaiki isi dan tampilan berdasarkan pemeriksaan};
    \draw[->] (arah) -- (sumber);
    \draw[->] (sumber) -- (isi);
    \draw[->] (isi) -- (keluaran);
    \draw[->] (keluaran) -- (perbaikan);
  \end{tikzpicture}
  \caption{Contoh diagram alur penyusunan naskah}
  \label{fig:contoh-alur}
\end{figure}

\subsection{Penyiapan Struktur}

Struktur awal disusun berdasarkan kebutuhan argumentasi. Setiap bab mempunyai tujuan yang dinyatakan pada pengantar, kemudian dibagi menjadi bagian yang mendukung tujuan tersebut. Judul dan isi diperiksa bersama agar suatu bagian tidak memiliki nama yang lebih luas daripada materinya. Penulis dapat membuat garis besar sebelum mengembangkan paragraf, lalu memperbaikinya setelah hubungan antargagasan lebih jelas.

\subsection{Penulisan dan Penandaan Unsur}

Setiap unsur yang dirujuk diberi label yang bermakna. Label untuk gambar, tabel, persamaan, dan bab dibedakan melalui awalan, sehingga penulis dapat mengenali jenis rujukan ketika menyunting sumber. Penggunaan label membantu mempertahankan hubungan narasi dengan unsur yang dirujuk ketika posisi unsur atau nomor bab berubah.

\subsubsection{Pemeriksaan Penanda}

Penanda perlu unik dalam satu dokumen. Label yang sama pada dua unsur dapat membuat rujukan menunjuk unsur yang berbeda dari maksud penulis. Nama label sebaiknya menggambarkan isi dan tidak hanya mengikuti nomor yang dapat berubah. Pada contoh ini, Gambar \ref{fig:contoh-alur} tetap dirujuk melalui pengenal yang menjelaskan alurnya.

\subsection{Pemeriksaan Keluaran}

Pemeriksaan keluaran dilakukan setelah isi dan rujukan disusun. Pembaca keluaran perlu dapat melihat judul, awal dan akhir tabel, keterangan gambar, rumus, serta daftar referensi dengan jelas. Penulis juga memeriksa apakah unsur yang panjang tetap terbaca ketika berpindah halaman. Kriteria pemeriksaan contoh ditata pada Tabel \ref{tab:contoh-metode}.

\begin{table}[htbp]
  \centering
  \caption{Contoh hubungan tahap dan bahan pemeriksaan}
  \label{tab:contoh-metode}
  \begin{tabular}{p{3.2cm}p{4.0cm}p{5.0cm}}
    \toprule
    Tahap & Bahan yang dibaca & Pertanyaan pemeriksaan \\
    \midrule
    Struktur & Judul bab dan bagian & Apakah setiap bagian mempunyai fungsi yang jelas? \\
    Sumber & Sitasi dan bibliografi & Apakah identitas sumber sesuai dengan dokumen asal? \\
    Unsur & Tabel, gambar, persamaan & Apakah narasi menjelaskan unsur yang dirujuk? \\
    Keluaran & Dokumen terkompilasi & Apakah isi dan tampilan terbaca secara utuh? \\
    Perbaikan & Catatan pemeriksaan & Apakah perubahan menyelesaikan persoalan yang dicatat? \\
    \bottomrule
  \end{tabular}
\end{table}

\section{Penyajian dan Analisis}

Keluaran disajikan berdasarkan pertanyaan yang hendak dijawab. Unsur dengan sifat sebanding dapat ditempatkan dalam tabel, sedangkan hubungan tahapan dapat diperlihatkan melalui diagram. Informasi yang membutuhkan uraian alasan tetap disampaikan melalui paragraf. Pemilihan bentuk penyajian bertujuan mengurangi beban membaca tanpa menghilangkan kondisi dan batas bukti.

Pada penelitian sesungguhnya, analisis perlu dijelaskan sebelum hasil ditafsirkan. Jika digunakan ukuran numerik, definisi, arah makna, satuan, serta aturan penanganan kondisi khusus perlu ditulis. Jika analisis berupa pembacaan tematik, penulis perlu menerangkan unit, kategori, dan cara mengikat kategori dengan bukti. Metode yang ditulis harus memungkinkan pembaca memahami asal setiap hasil yang dilaporkan.

\section{Ketertelusuran}

Sumber naskah dan bahan pendukung disimpan dengan nama yang menunjukkan fungsi. Perubahan penting dicatat agar alasan perbaikan dapat ditemukan kembali. Pada template, pemisahan berkas per bab membantu pengguna menyunting bagian tertentu sambil mempertahankan satu berkas utama untuk penyusunan keseluruhan dokumen.

Penelitian dengan bahan yang tidak dapat dibagikan tetap memerlukan penjelasan tentang keberadaan dan batas akses bahan tersebut. Catatan ketertelusuran dapat menyertakan pengenal serta prosedur pemeriksaan tanpa mempublikasikan isi yang dibatasi. Lampiran \ref{lamp:contoh} memperlihatkan cara menata catatan pendukung secara umum; pengguna perlu menggantinya dengan informasi yang sesuai dengan penelitiannya.

\chapter{HASIL}
\label{bab:hasil}

\textbf{Bab ini merupakan contoh format pelaporan dan tidak memuat hasil penelitian riil.} Tidak ada skor, jumlah sampel empiris, keberhasilan eksperimen, ataupun perbandingan metode yang dinyatakan telah diperoleh. Tabel berisi nama komponen dan deskripsi keluaran yang perlu disiapkan oleh pengguna. Ganti seluruh materi contoh dengan hasil yang dapat ditelusuri ke data dan prosedur penelitian sendiri.

\section{Cakupan Keluaran}

Pengantar hasil perlu menjelaskan keluaran yang tersedia, hubungan dengan metodologi, dan batas penyajiannya. Jika suatu prosedur tidak menghasilkan keluaran yang direncanakan, kondisi tersebut perlu dinyatakan. Tidak tersedianya hasil adalah informasi yang berbeda dari hasil bernilai nol. Memisahkan keduanya membantu pembaca menilai cakupan jawaban tanpa menganggap seluruh tahap selalu berhasil.

Tabel \ref{tab:contoh-keluaran} menunjukkan format ringkas untuk mengelompokkan komponen. Nama pada tabel merupakan unsur dokumen, bukan daftar temuan penelitian. Bentuk ini dapat diganti dengan komponen yang sesuai dengan bidang pengguna, misalnya kelompok data, tahap analisis, kategori pengamatan, atau keluaran pengukuran.

\begin{table}[htbp]
  \centering
  \caption{Contoh format ringkasan komponen keluaran, tanpa hasil empiris}
  \label{tab:contoh-keluaran}
  \begin{tabular}{p{4.0cm}p{8.6cm}}
    \toprule
    Nama komponen & Deskripsi keluaran \\
    \midrule
    Identitas bahan & Keterangan asal, cakupan, dan batas penggunaan bahan \\
    Prosedur analisis & Keterangan hubungan langkah dengan keluaran yang disajikan \\
    Tabel ringkasan & Susunan informasi yang dapat dibandingkan dalam satu tampilan \\
    Gambar alur & Representasi hubungan tahapan yang dijelaskan dalam narasi \\
    Catatan pemeriksaan & Penjelasan kondisi yang perlu diperbaiki atau diperiksa kembali \\
    \bottomrule
  \end{tabular}
\end{table}

\section{Penyajian Unsur Dokumen}

\subsection{Tabel dan Penjelasan}

Tabel membutuhkan judul yang menjelaskan isi tanpa bergantung sepenuhnya pada paragraf lain. Judul tetap ringkas, sedangkan aturan pembacaan yang lebih panjang dapat diterangkan dalam narasi. Nama kolom harus menunjukkan jenis informasi yang disajikan. Jika tabel penelitian memuat besaran, satuan dan arti nilai kosong perlu disebutkan dengan cara yang konsisten.

Sesudah tabel, penulis dapat menjelaskan hubungan yang paling relevan dengan pertanyaan penelitian. Uraian tidak perlu mengulang seluruh isi sel dalam bentuk paragraf. Penjelasan dipilih berdasarkan peran tabel dalam argumentasi. Dalam contoh ini, Tabel \ref{tab:contoh-keluaran} dipakai untuk menunjukkan bahwa komponen dan deskripsinya mempunyai kolom yang terpisah serta dapat dirujuk dari narasi.

\subsection{Gambar dan Persamaan}

Gambar \ref{fig:contoh-alur} pada Bab \ref{bab:metodologi} memperlihatkan rujukan silang antarbab. Penulis dapat merujuk gambar yang telah diterangkan sebelumnya apabila gambarnya masih relevan. Pengulangan gambar perlu dipertimbangkan berdasarkan kebutuhan pembaca dan panjang pembahasan, bukan dilakukan hanya karena bab berubah.

Persamaan \ref{eq:contoh-rerata} dapat dirujuk untuk menjelaskan definisi suatu besaran. Nilai yang dihasilkan pada penelitian sesungguhnya harus mempunyai satuan, definisi, dan asal penghitungan yang dapat diperiksa. Contoh ini tidak memasukkan nilai ke persamaan karena tidak menyediakan observasi empiris. Penyajian formula dan penyajian hasil pengukuran memiliki peran berbeda.

\section{Contoh Tabel yang Melintasi Halaman}

Tabel \ref{tab:contoh-pemeriksaan-panjang} menggunakan lingkungan tabel panjang agar dapat berpindah halaman dengan judul kolom yang diulang. Keterangan ``(lanjutan)'' membantu pembaca mengenali hubungan bagian tabel pada halaman berikutnya. Kolom pertama menamai komponen dan kolom kedua menjelaskan bahan yang dibutuhkan. Isi tetap berupa petunjuk penulisan, tanpa status pemeriksaan yang diklaim sudah lulus.

\begingroup
\setlength{\tabcolsep}{5pt}
\begin{longtable}{p{3.5cm}p{9.0cm}}
  \caption{Contoh daftar komponen dan deskripsi keluaran untuk tabel panjang}
  \label{tab:contoh-pemeriksaan-panjang}\\
  \toprule
  Nama komponen & Deskripsi keluaran \\
  \midrule
  \endfirsthead
  \caption[]{Contoh daftar komponen dan deskripsi keluaran untuk tabel panjang (lanjutan)}\\
  \toprule
  Nama komponen & Deskripsi keluaran \\
  \midrule
  \endhead
  \midrule
  \multicolumn{2}{r}{Bersambung ke halaman berikutnya} \\
  \endfoot
  \bottomrule
  \endlastfoot
  Judul naskah & Judul yang menjelaskan fokus penelitian dan sesuai dengan judul yang telah disepakati. Penggantian judul dilakukan melalui metadata agar bagian yang memakainya tetap konsisten. \\
  Identitas penulis & Nama dan nomor identitas yang diperiksa terhadap data akademik. Nilai pada contoh merupakan penanda generik yang perlu diganti sebelum naskah digunakan. \\
  Identitas pembimbing & Nama pembimbing sesuai penetapan yang berlaku. Nomor identitas serta pengesahan dicantumkan hanya berdasarkan informasi yang dapat dipertanggungjawabkan. \\
  Persetujuan & Redaksi formulir, ruang tanda tangan, dan tanggal berdasarkan proses persetujuan sebenarnya. Halaman contoh tidak menyatakan bahwa persetujuan telah diberikan. \\
  Orisinalitas & Redaksi pernyataan sesuai ketentuan yang berlaku dan praktik penulisan sebenarnya. Ruang formulir tetap kosong sampai penulis menyiapkan pernyataan yang dapat dipertanggungjawabkan. \\
  Prakata & Uraian singkat mengenai konteks penyusunan dan bantuan yang benar-benar diterima. Nama atau ucapan terima kasih tidak ditambahkan secara fiktif. \\
  Abstrak & Ringkasan masalah, tujuan, metode, hasil, dan batas kesimpulan pada penelitian sesungguhnya. Abstrak contoh menerangkan fungsi template karena hasil empiris tidak disediakan. \\
  Abstract & Ringkasan berbahasa Inggris yang mempunyai makna setara dengan abstrak Indonesia. Istilah, cakupan, dan batas hasil perlu diperiksa pada kedua bahasa. \\
  Kata kunci & Istilah yang membantu pembaca mengenali bidang dan fokus naskah. Istilah dipilih berdasarkan materi sebenarnya dan mengikuti ketentuan jumlah yang berlaku. \\
  Daftar isi & Judul bab dan bagian yang berasal dari struktur sumber. Setelah isi berubah, keluaran diperiksa agar nomor halaman serta hierarki judul mengikuti susunan terbaru. \\
  Daftar tabel & Judul tabel yang konsisten dengan tabel di dalam naskah. Judul yang panjang dapat ditata agar daftar tetap terbaca tanpa menghilangkan makna utamanya. \\
  Daftar gambar & Judul gambar dan lokasi halaman yang sesuai dengan penyajian dalam isi. Gambar perlu dirujuk dan dijelaskan, termasuk apabila berasal dari sumber lain. \\
  Daftar lampiran & Pengenal materi pendukung yang dapat ditemukan kembali oleh pembaca. Materi yang dirujuk sebagai lampiran harus benar-benar terdapat pada bagian akhir naskah. \\
  Pendahuluan & Hubungan masalah, tujuan, manfaat, serta batasan yang menjadi dasar penelitian. Uraian disusun berdasarkan sumber dan kondisi yang dapat diperiksa. \\
  Landasan kepustakaan & Konsep dan penelitian terdahulu yang digunakan dalam argumentasi. Setiap sumber memiliki peran yang dapat dikenali dalam pemilihan metode atau penafsiran hasil. \\
  Metodologi & Prosedur, unit pengamatan, bahan, dan aturan analisis yang menjawab rumusan masalah. Bagian ini menjelaskan asal keluaran yang dilaporkan pada bab hasil. \\
  Hasil & Bukti yang benar-benar diperoleh beserta cakupan dan kondisi pelaporannya. Prosedur yang tidak menghasilkan keluaran dijelaskan tanpa membuat nilai pengganti. \\
  Pembahasan & Hubungan antara bukti, pertanyaan penelitian, konsep, dan keterbatasan. Penafsiran dibatasi oleh cakupan bahan serta prosedur yang digunakan. \\
  Penutup & Jawaban ringkas yang mengikuti bukti dan pembahasan. Saran berhubungan dengan keterbatasan yang telah dijelaskan, bukan tambahan hasil yang belum dianalisis. \\
  Sitasi & Pengenal sumber yang mendukung pernyataan terkait. Penulis memeriksa bahwa kutipan dan parafrasa sesuai dengan isi serta batas sumber asal. \\
  Daftar referensi & Identitas bibliografis sumber yang dipakai dalam naskah. Jenis sumber, penulis, judul, dan alamat diperiksa tanpa mengisi data penerbitan dengan nilai rekaan. \\
  Persamaan & Definisi simbol, syarat penggunaan, serta rujukan pada narasi. Persamaan yang memerlukan beberapa baris ditata dengan nomor yang sesuai kebutuhan rujukan. \\
  Lampiran & Materi pendukung yang terlalu rinci untuk alur utama tetapi diperlukan untuk pemeriksaan. Bahan yang dibatasi aksesnya tidak disertakan tanpa dasar izin yang sesuai. \\
  Keluaran akhir & Dokumen yang diperiksa kembali sesudah perubahan terakhir. Pemeriksaan mencakup tampilan, rujukan, identitas, serta kelengkapan yang diperlukan untuk penggunaan akademik. \\
\end{longtable}
\endgroup

\section{Kelengkapan Pelaporan}

Pelaporan perlu mempertahankan hubungan antara bahan, prosedur, dan keluaran. Apabila informasi disingkat dalam tabel utama, rincian dapat diletakkan pada lampiran dan dirujuk secara jelas. Informasi penting tentang batas bukti tidak sebaiknya hanya diletakkan di berkas terpisah yang tidak disebut dalam naskah.

Pada contoh ini, Lampiran \ref{lamp:contoh} menyediakan susunan catatan pendukung yang generik. Pengguna dapat mengganti catatan tersebut dengan bahan yang relevan dan berizin. Bab \ref{bab:pembahasan} memperlihatkan cara meneruskan pelaporan menuju penafsiran, tanpa menjadikan contoh format sebagai klaim bahwa hasil ilmiah telah tersedia.

\chapter{PEMBAHASAN}
\label{bab:pembahasan}

Bab ini menunjukkan susunan pembahasan dan cara menghubungkan bukti dengan argumentasi. Materi tetap berupa petunjuk penulisan. Karena contoh tidak memuat hasil empiris, bab ini tidak menyimpulkan kinerja metode, keberhasilan sistem, ataupun pengaruh suatu perlakuan.

\section{Hubungan Keluaran dengan Pertanyaan}

Pembahasan dimulai dengan pertanyaan yang hendak dijawab, kemudian menunjuk keluaran yang menjadi dasar penafsiran. Hubungan tersebut membantu pembaca melihat bagaimana hasil menjawab masalah. Penulis dapat menyusun pembahasan per pertanyaan, per tema, atau per kelompok hasil, selama alurnya dijelaskan dan tidak memisahkan informasi yang perlu dibaca bersama.

Dalam contoh, Tabel \ref{tab:contoh-keluaran} memisahkan nama komponen dari deskripsinya. Pemisahan ini memperlihatkan cara merujuk unsur yang menjadi bahan penjelasan. Pada skripsi sesungguhnya, penulis perlu mengganti komponen itu dengan bukti yang relevan. Menunjuk nomor tabel tidak cukup; pembahasan perlu menjelaskan hubungan yang ditemukan dan mengapa hubungan tersebut menjawab pertanyaan.

\subsection{Bukti dan Penafsiran}

Bukti menyatakan apa yang diamati atau dihitung, sedangkan penafsiran menjelaskan makna bukti dalam konteks penelitian. Keduanya perlu dibedakan agar pembaca dapat menilai apakah penjelasan mengikuti keluaran. Suatu perbedaan yang diamati tidak otomatis membuktikan alasan terjadinya perbedaan. Penjelasan alternatif dapat disampaikan dengan menunjukkan bagian yang masih membutuhkan pemeriksaan.

Pembahasan yang kuat dapat dimulai dari temuan yang paling relevan, dilanjutkan dengan bukti pendukung, kemudian menghubungkannya dengan konsep. Urutan ini tidak harus mengikuti urutan pelaksanaan pekerjaan. Penulis memilih alur yang memudahkan pembaca menilai kesimpulan, sekaligus mempertahankan kondisi ketika bukti diperoleh.

\subsection{Hubungan dengan Landasan}

Landasan dalam Bab \ref{bab:landasan} berperan ketika penulis membandingkan atau menjelaskan hasil. Konsep tidak hanya diulang sebagai definisi, melainkan digunakan untuk menerangkan hubungan tertentu. Bila hasil berbeda dari laporan terdahulu, penulis perlu memeriksa perbedaan objek, cakupan, prosedur, dan ukuran sebelum menyatakan bahwa satu hasil bertentangan dengan hasil lain.

Pada contoh ini, pemisahan struktur dari tampilan digunakan untuk menerangkan kegunaan label dan rujukan silang. Hubungan tersebut relevan dengan penataan naskah. Pada penelitian pengguna, konsep dokumentasi ini diganti dengan landasan bidang yang sebenarnya. Pembahasan tidak sebaiknya mempertahankan uraian template ketika pertanyaan penelitiannya sudah berubah.

\section{Batas Penafsiran}

Tabel \ref{tab:contoh-keterbatasan} memperlihatkan format pembahasan keterbatasan. Kolom pertama menamai batas, kolom kedua menjelaskan akibatnya bagi pembacaan, dan kolom ketiga menempatkan tindak lanjut. Isi tabel menjelaskan batas contoh naskah; tabel tidak mencatat keterbatasan suatu penelitian yang belum dilaksanakan.

\begin{table}[htbp]
  \centering
  \caption{Contoh penataan batas dan tindak lanjut naskah}
  \label{tab:contoh-keterbatasan}
  \begin{tabular}{p{3.0cm}p{4.4cm}p{4.8cm}}
    \toprule
    Batas contoh & Akibat bagi pembacaan & Tindak lanjut pengguna \\
    \midrule
    Tidak ada data empiris & Format tidak dapat dibaca sebagai temuan & Ganti isi dengan bukti penelitian sendiri \\
    Landasan dokumentasi & Sumber belum membahas bidang penelitian pengguna & Susun tinjauan yang relevan dan terverifikasi \\
    Formulir generik & Halaman awal belum menjadi dokumen persetujuan & Sesuaikan dengan formulir serta persetujuan aktual \\
    Susunan enam bab & Hierarki perlu disesuaikan dengan jenis penelitian & Periksa pedoman dan arahan pembimbing \\
    \bottomrule
  \end{tabular}
\end{table}

Penjelasan keterbatasan perlu menunjukkan akibat yang konkret bagi kesimpulan. Pernyataan umum bahwa suatu penelitian memiliki keterbatasan belum membantu pembaca menentukan bagian mana yang dapat digunakan. Penulis sebaiknya menyebut cakupan yang belum diperiksa, alasan pembatasan, dan jenis klaim yang masih membutuhkan bukti tambahan.

\section{Implikasi bagi Penulisan}

Rujukan silang membantu menghubungkan alur naskah, tetapi pengguna tetap perlu membaca ulang paragraf setelah unsur dipindahkan. Nomor yang benar tidak memastikan bahwa kalimat masih menjelaskan unsur yang dimaksud. Demikian pula, sitasi yang tampil tidak memastikan bahwa sumber mendukung pernyataan. Pemeriksaan isi dan pemeriksaan tampilan perlu dilakukan bersama.

Materi pendukung ditempatkan menurut kebutuhan pembaca. Rincian yang terlalu panjang dapat berada pada lampiran, tetapi hasil yang menentukan kesimpulan perlu tetap dijelaskan pada alur utama. Dalam contoh, tabel panjang memperlihatkan cara melanjutkan penyajian ke halaman berikutnya. Pada penelitian sesungguhnya, pemilihan panjang dan letak tabel mengikuti kegunaan informasinya.

\section{Persiapan Penutup}

Sebelum menyusun penutup, penulis dapat memeriksa apakah setiap rumusan masalah telah mempunyai jawaban dengan dasar yang disebutkan. Jawaban yang masih bergantung pada asumsi harus mempertahankan asumsi tersebut. Jika suatu pertanyaan belum dapat dijawab, keadaan itu perlu dinyatakan dan dapat menjadi dasar saran yang terarah.

Bab \ref{bab:penutup} menyediakan contoh susunan simpulan dan saran. Pada naskah akhir, penutup merupakan hasil pemadatan argumentasi yang telah dibahas, sehingga tidak menambahkan bukti baru. Pengguna perlu memastikan bahwa istilah dan batas jawaban tetap sama ketika uraian panjang dipadatkan menjadi simpulan.

\chapter{PENUTUP}
\label{bab:penutup}

Bab ini menunjukkan penempatan simpulan dan saran dalam susunan contoh. Uraiannya merangkum fungsi template dan tidak menyatakan adanya hasil penelitian ilmiah. Ketika template dipakai untuk skripsi, ganti bagian ini dengan jawaban atas rumusan masalah berdasarkan hasil dan pembahasan yang sudah tersedia.

\section{Simpulan}

Contoh naskah menyediakan struktur enam bab, halaman awal, daftar referensi, dan lampiran. Unsur tabel, diagram vektor, persamaan, serta sitasi ditampilkan untuk memperlihatkan pola penulisan dan rujukan silang. Identitas, isi ilmiah, redaksi administratif, dan sumber yang relevan perlu disiapkan oleh pengguna sesuai penelitian serta ketentuan yang berlaku.

Tabel \ref{tab:contoh-rangkuman} memperlihatkan cara memetakan fokus dengan bentuk jawaban. Isi tabel adalah rangkuman fungsi contoh, bukan kesimpulan yang diperoleh dari pengujian. Penulis dapat memakai susunannya untuk memeriksa apakah rumusan masalah dan simpulan pada skripsi sesungguhnya mempunyai hubungan yang jelas.

\begin{table}[htbp]
  \centering
  \caption{Contoh pemetaan fokus penulisan dan bentuk rangkuman}
  \label{tab:contoh-rangkuman}
  \begin{tabular}{p{3.8cm}p{8.8cm}}
    \toprule
    Fokus penulisan & Bentuk rangkuman pada contoh \\
    \midrule
    Hubungan antarbagian & Bab dihubungkan melalui pengantar, sistematika, dan rujukan silang \\
    Unsur naskah & Tabel, gambar, persamaan, sitasi, serta lampiran mempunyai contoh yang dapat disunting \\
    Pemeriksaan dokumen & Isi dan tampilan perlu diperiksa setelah materi contoh diganti \\
    \bottomrule
  \end{tabular}
\end{table}

Dalam penelitian sesungguhnya, simpulan sebaiknya menyatakan jawaban secara langsung dan menjaga batas bukti. Pernyataan yang lebih luas daripada objek, sumber, atau kondisi yang diperiksa dapat memberi makna yang tidak didukung. Jumlah simpulan disesuaikan dengan pertanyaan serta hubungan antartemuan, sehingga bukan sekadar mengikuti jumlah poin pada contoh.

\section{Saran}

Pengguna disarankan menyiapkan garis besar isi sebelum mengganti paragraf contoh. Langkah ini membantu menyesuaikan fungsi tiap bab dengan jenis penelitian. Setelah materi diganti, rujukan silang, judul tabel, daftar referensi, dan lampiran diperiksa kembali agar tetap sesuai dengan alur serta unsur yang benar-benar tersedia.

Redaksi halaman persetujuan dan orisinalitas perlu disesuaikan dengan formulir yang berlaku. Nama dan nomor identitas harus diperiksa sebelum berkas dipakai untuk keperluan akademik. Ruang tanggal dan tanda tangan hanya diisi berdasarkan proses sebenarnya. Penggunaan template tidak menggantikan persetujuan pembimbing terhadap naskah.

Pada skripsi akhir, saran diturunkan dari keterbatasan yang telah dibahas. Saran dapat menyebut pemeriksaan tambahan, sumber yang perlu diperluas, atau keputusan metode yang perlu diuji pada kondisi lain. Penulis perlu menjelaskan alasan usulan tersebut tanpa menjadikannya sebagai janji keberhasilan penelitian lanjutan.

\printbibliography[heading=bibintoc,title={DAFTAR REFERENSI}]
\appendix
\skripsiappendix{Contoh Materi Pendukung}
\label{lamp:contoh}

Lampiran ini memperlihatkan penempatan informasi pendukung di luar alur utama. Semua isinya merupakan petunjuk format, tanpa data penelitian pribadi. Pengguna dapat menggantinya dengan instrumen, rincian prosedur, daftar parameter, hasil pemeriksaan, atau bahan pendukung lain yang diperlukan pembaca dan diizinkan untuk dibagikan.

\section*{Keterangan Materi}

Setiap bahan pendukung perlu memiliki fungsi yang jelas. Jelaskan bagian naskah yang memerlukan bahan tersebut, asal bahan, serta batas penggunaannya. Apabila lampiran memuat salinan dokumen pihak lain, izin dan atribusi diperiksa sebelum salinan disertakan. Lampiran tidak sebaiknya menjadi tempat memasukkan seluruh berkas pengerjaan tanpa menilai kegunaannya bagi pembaca.

Tabel \ref{tab:contoh-lampiran} menunjukkan format catatan yang dapat diperiksa. Kolom tidak menyatakan bahwa suatu pemeriksaan sudah selesai. Pengguna dapat menambahkan status aktual serta tanggal hanya setelah ada bukti yang sesuai.

\begin{table}[htbp]
  \centering
  \caption{Contoh keterangan bahan pendukung}
  \label{tab:contoh-lampiran}
  \begin{tabular}{p{4.0cm}p{8.6cm}}
    \toprule
    Nama bahan & Keterangan yang disiapkan \\
    \midrule
    Sumber naskah & Nama berkas dan bagian yang menyimpan materi terkait \\
    Data bibliografis & Pengenal entri, sumber identitas, dan tanggal pemeriksaan \\
    Instrumen & Asal, tujuan, cara penggunaan, dan izin yang diperlukan \\
    Catatan perubahan & Bagian yang diperbaiki, alasan, dan dasar pemeriksaan \\
    Keluaran dokumen & Lingkungan penyusunan dan pemeriksaan tampilan terakhir \\
    \bottomrule
  \end{tabular}
\end{table}

\section*{Contoh Catatan Rujukan Silang}

Pola rujukan silang pada template dapat dipelajari melalui Tabel \ref{tab:contoh-keluaran}, Gambar \ref{fig:contoh-alur}, dan Persamaan \ref{eq:contoh-rerata}. Ketiganya mempunyai label unik yang mengikuti jenis unsur. Bab \ref{bab:hasil} juga merujuk lampiran ini, sehingga hubungan dari isi utama ke materi pendukung dan kembali ke isi utama dapat diperiksa.

Pengguna sebaiknya memberi nama label yang menggambarkan isi, lalu memanggil label ketika menyebut unsur dalam paragraf. Hindari memasukkan nomor tabel atau gambar secara manual karena nomor dapat berubah setelah penyuntingan. Nomor yang dibentuk otomatis tetap perlu diperiksa bersama makna kalimat yang menggunakannya.

\section*{Catatan sebelum Mengganti Contoh}

Simpan satu salinan sumber contoh untuk mempelajari pola penulisan. Pada naskah kerja, ganti identitas, abstrak, prakata, isi bab, dan bibliografi dengan materi sebenarnya. Hapus gambar atau tabel contoh apabila tidak lagi digunakan, lalu periksa rujukan yang sebelumnya menunjuk unsur tersebut. Materi yang dihapus tidak boleh meninggalkan nomor rujukan yang belum terdefinisi.

Penulis tetap bertanggung jawab terhadap isi ilmiah dan proses persetujuan naskah. Perangkat penyusunan membantu mengatur tampilan dan hubungan penanda, tetapi tidak memverifikasi kebenaran suatu temuan. Pemeriksaan akhir mencakup isi, sumber, tampilan, serta ketentuan akademik yang berlaku pada saat naskah diserahkan.

\skripsiappendix{Contoh Lampiran Tambahan}
Lampiran kedua memperlihatkan bahwa penomoran berubah menjadi B. Judul pada daftar isi dan daftar lampiran menggunakan label LAMPIRAN, sedangkan tabel di dalam lampiran memakai awalan B sebagaimana Tabel \ref{tab:lampiran-b}.
\begin{table}[htbp]
\centering
\caption{Contoh format daftar berkas pendukung}
\label{tab:lampiran-b}
\begin{tabular}{ll}
\toprule
Jenis berkas & Keterangan \\
\midrule
Dokumen & Instrumen atau prosedur penelitian \\
Catatan & Ringkasan hasil pemeriksaan \\
\bottomrule
\end{tabular}
\end{table}

\end{document}
